\section{LLVM IR 与优化}
\label{sec:ir}

\subsection{控制流与内存效应的显式化}

LLVM IR 采用带类型的 SSA 表示，函数由基本块组成，每个基本块以控制流终结指令结束\cite{llvm-langref-2026}。在 \code{-O0} 产物中，可寻址局部变量通常对应 \code{alloca} 栈对象，初始化和赋值形成 \code{store}，表达式取值形成 \code{load}。源级 \code{while} 被降低为条件、循环体和退出基本块，循环体至条件块的分支构成回边。此处的 \code{alloca} 表示抽象栈对象，具体栈偏移仍由目标代码生成决定。

参考 LLVM IR 直接以 SSA 形式表示阶乘递推。清单~\ref{lst:phi} 中两个 $\phi$ 操作依据实际执行的前驱边选择输入定义：首次迭代从 \code{entry} 取得 2 和 1，后续迭代从 \code{body} 取得更新值。支配关系和支配边界为一般 SSA 构造中的合流位置提供理论基础\cite{cytron-et-al-ssa-1991}。案例的循环回边要求 \code{factor} 与 \code{result} 在循环头合并初始值和迭代值。

\begin{center}
\begin{minipage}{0.92\linewidth}
\begin{lstlisting}[caption={阶乘循环在参考 LLVM IR 中的 SSA 表示。},label={lst:phi}]
loop:
  %factor = phi i32 [ 2, %entry ], [ %next_factor, %body ]
  %result = phi i32 [ 1, %entry ], [ %next_result, %body ]
  %continue = icmp sle i32 %factor, %n
  br i1 %continue, label %body, label %done
body:
  %next_result = mul i32 %result, %factor
  %next_factor = add i32 %factor, 1
  br label %loop
\end{lstlisting}
\end{minipage}
\end{center}

函数调用在 IR 中同时编码返回类型、目标符号和参数类型。例如 \code{call i32 @factorial(i32 \%bounded)} 返回 32 位整数，\code{putint} 与 \code{putch} 则声明为 \code{void}。夹取条件分别采用 \code{icmp slt} 和 \code{icmp sgt}，循环条件采用 \code{icmp sle}，明确表达有符号整数比较。参考 IR 的 \code{mul i32} 未附加 \code{nsw} 标志，因而没有引入额外的无溢出假设；已声明输入路径上的乘积由上界保证有效。

\subsection{SSA 构造与标量优化}

LLVM 优化流水线由分析和变换通道组合，其默认构成随版本、目标和优化级别演化\cite{llvm-passes-2026,llvm-newpm-2026}。案例中的局部标量地址不逃逸，使标量提升能够把栈对象的读写改写为 SSA 定义和使用；常量传播、控制流简化及内联进一步改变基本块和调用边界。由此产生的文本结构差异以语义保持的程序变换为前提。

表~\ref{tab:structure} 合并报告 IR 与汇编层的结构测量。RV64 目标下，IR 指令数由 56 降至 30，\code{alloca/load/store} 从 9/13/13 降至 0/0/0，反映出局部内存状态的提升与整体结构收缩。汇编近似指令行由 77 降至 45，而助记符种类由 13 增至 15，表明文本规模和操作种类并不单调对应。这些静态指标用于刻画表示结构，性能解释不在其测量范围内。

\begin{table}[tb]
\caption{Clang 18.1.3 在 RV64 目标上的表示结构测量。}
\label{tab:structure}
\centering\small
\begin{tabular}{@{}lrrrrrr@{}}
\toprule
级别 & IR 指令 & \code{alloca} & \code{load} & \code{store} & 汇编行 & 助记符\\
\midrule
\code{-O0} & 56 & 9 & 13 & 13 & 77 & 13\\
\code{-O2} & 30 & 0 & 0 & 0 & 45 & 15\\
\bottomrule
\end{tabular}
\end{table}

优化后的 IR 仍使用不限数量的 SSA 名称和目标相对独立的整数操作。目标后端接收其数据依赖、控制流和内存语义，并将这些信息映射到有限物理寄存器、具体指令和调用约定。
